\documentclass[10pt,a4paper]{amsart}
\usepackage[margin=19mm]{geometry}
\usepackage{amsmath,amssymb,mathtools}
\usepackage[scr=rsfs,cal=euler]{mathalfa}
\usepackage{xfrac}
\usepackage[unicode,hidelinks,bookmarks=false]{hyperref}
\newcommand{\ie}{i.e.\ }
\newcommand{\cf}{cf.\ }
\newcommand{\resp}{respectively\ }
\newcommand{\Subsection}{\subsection}
\providecommand{\nobreakdash}{\nobreak\hbox{-}}
\setlength{\emergencystretch}{2em}
\multlinegap0pt
\usepackage{bm}
\usepackage{bbm}
\usepackage{dsfont}
\usepackage{xspace}
\usepackage{cancel}
\usepackage{mathtools}
\usepackage{amsthm}
\makeatletter
\@ifundefined{theorem}{\newtheorem{theorem}{Theorem}}{}
\makeatother
\title[Is the Frequency Principle always valid?]{Is the Frequency Principle always valid?}
\author{Qijia Zhai}
\date{Source: arXiv:2508.17323v1 (CC BY 4.0)}
\begin{document}
\maketitle

\noindent\emph{Source and licence.} Is the Frequency Principle always valid?. Version arXiv:2508.17323v1 (CC BY 4.0), distributed under the Creative Commons Attribution 4.0 International licence, as recorded in the arXiv licence metadata for this exact version and shown on the abstract page https://arxiv.org/abs/2508.17323v1. Full rights notice: licenses/arxiv-2508.17323-CC-BY-4.0.txt.\par\smallskip

\noindent\textit{Editorial scope.} Counted unit: the Theorem whose source label is \texttt{thm:FPviolationSphere} in the article (source0.tex lines 711--720, source label \texttt{thm:FPviolationSphere}), delivered together with the article's own complete original proof of it (source0.tex lines 721--769, one \texttt{proof} environment of 49 lines). No mathematics is written, repaired or re-derived: statement and proof are the article's own text line for line. Required context retained uncounted: neuralnetwork-sphere (lines 151--159); errorxi-sphere (lines 198--204); D\_ellfixed (lines 669--673). Every definition, display and label of those lines is retained unchanged. Omitted from this package and not redistributed: 1--150, 160--197, 205--668, 674--710, 770--2225. Nothing inside a retained excerpt is omitted or altered: every retained body is delivered exactly as the article's lines are written, so no omission receipt of any kind accompanies this package. Citations: the retained excerpts carry no citation command at all, so no bibliography entry of the article is transcribed and no reference is redistributed. Reference adaptation: none was needed and none was made; every reference inside the delivered unit resolves inside it, so no unresolved reference or citation is produced. Printed slips kept literal: no printed word, symbol, hypothesis or number of the counted statement or of its proof was altered. Layout and class: the article's own class is \texttt{siamart250106} with packages including amsmath, amsfonts, amsopn, amssymb, graphicx, multirow, relsize, makecell, subfigure, url, enumitem, algorithm, algorithmic, ulem; the wrapper is the lane's pinned \texttt{amsart} editorial wrapper at 10pt on A4, which restates the theorem-like environments the retained text needs and adds the article's own macro definitions that the retained text needs (none), with extra wrapper packages (bm, bbm, dsfont, xspace, cancel, mathtools). The delivered numbering is the unit's own. Assets: no retained span contains a figure environment, a graphics-inclusion command, a PDF-page inclusion, a table or a TikZ picture; no image file is copied into this package, the package declares no asset dependency and no image file is redistributed with it. Licence: version arXiv:2508.17323v1 of this article is distributed under the Creative Commons Attribution 4.0 International licence, as recorded in the arXiv licence metadata for this exact version and shown on the abstract page https://arxiv.org/abs/2508.17323v1. A full rights notice is delivered at licenses/arxiv-2508.17323-CC-BY-4.0.txt. Characters: every retained excerpt is pure ASCII, so the delivered engine, XeLaTeX with its default Latin Modern text font, reports no missing character.
\begin{equation}\label{neuralnetwork-sphere}
u\bigl(\tau,\phi;\theta\bigr)
\;=\;
\sum_{i=1}^m 
a_i\,\sigma\bigl(w_i^\top x(\tau,\phi)\bigr)
\;=\;
\sum_{i=1}^m 
a_i\,\mathrm{ReLU}\bigl(w_i^\top x(\tau,\phi)\bigr).
\end{equation}

\begin{equation}\label{errorxi-sphere}
    D(\tau,\phi)
\;=\;
u\bigl(\tau,\phi;\theta\bigr)
\;-\;
h(\tau,\phi),
\end{equation}

\begin{equation}\label{D_ellfixed}
D_{\ell}
\;:=\;
m\Bigl(\tfrac{2\pi}{3}\!\!\sum_{k=1}^m a_k - C(h)\Bigr)\,c_{\ell}.
\end{equation}

\begin{theorem}[General Conditions for FP Violation on the Sphere]
\label{thm:FPviolationSphere}
Let \(\{w_i\}_{i=1}^m \subset S^2\) be fixed unit vectors and align along the polar axis, and consider the shallow ReLU network \eqref{neuralnetwork-sphere}. Let \(D(\tau,\phi)\) defined in \eqref{errorxi-sphere} be the error function on the sphere.  
Suppose that at any training time $t$, if $$
\sum_{k=1}^m a_k < \frac{3}{2\pi}C(h)$$
with $C(h) = \int_{0}^{2\pi}\!\!\int_{0}^{\tfrac{\pi}{2}}
h(\tau',\phi')\,\cos(\tau')\,\sin(\tau')\,d\tau'\,d\phi'$, then for all \(\ell\), the frequency principle can be \emph{violated}.  
In other words, the exponential decay of \(c_{\ell}\) can be surmounted by such error terms at sufficiently low 
frequencies, enabling high-frequency components to persist.
\end{theorem}

\begin{proof}[Proof of Theorem \ref{thm:FPviolationSphere}]
    This proof can be directly calculated from \eqref{D_ellfixed} that $c_{\ell}$ is positive for all $\ell$ and $\Bigl(\tfrac{2\pi}{3}\!\!\sum_{k=1}^m a_k - C(h)\Bigr)$ is negative. So $\frac{\partial D_{\ell}}{\partial t} > 0$, which means the neural network is not convergent.
% We claim that for $\ell \le 1$, the ratio of consecutive spherical-harmonic coefficients 
% $\bigl|c_{\ell+1}\bigr|/\bigl|c_{\ell}\bigr|$ can exceed 1, thereby contradicting the standard frequency principle.  

% \noindent
% Recall that each single-neuron ReLU function 
% \(\mathrm{ReLU}\!\bigl(w_i^\top x(\tau,\phi)\bigr)\) (with \(w_i\) aligned along the polar axis) admits an expansion 
% \(\sum_{\ell=0}^\infty c_{\ell}\,Y_{\ell}^0(\tau,\phi)\), whose coefficients for \(\ell\ge0\) satisfy
% \[
% \frac{\bigl|c_{\ell+1}\bigr|}{\bigl|c_{\ell}\bigr|}
% \;=\;
% \frac{\sqrt{2\ell + 3}}{2\,\sqrt{\,2\ell + 1\,}}
% \,\cdot\,
% \frac{(\ell + 1)^2 + 3(\ell + 1) + 2}{\ell^2 + 3\ell + 2}.
% \]
% For \(\ell=0\), direct computation of \(c_0^{(i)}\) and \(c_1^{(i)}\) gives
% \[
% \bigl|c_0^{(i)}\bigr| \;=\; \frac{\sqrt{\pi}}{2},
% \quad
% \bigl|c_1^{(i)}\bigr| \;=\; \frac{2\sqrt{3\pi}}{3}.
% \]
% Thus
% \[
% \frac{\bigl|c_1^{(i)}\bigr|}{\bigl|c_0^{(i)}\bigr|}
% \;=\;
% \frac{4\,\sqrt{3}}{3}
% \,\approx\,
% 2.31
% \,>\,
% 1.
% \]


% By hypothesis, $\sum_{k=1}^m a_k > \frac{3}{2\pi}C(h)$. Hence, in the update term for the $\ell$-th mode, the factors $c_{\ell}$ all combine \emph{without sign cancellations}.  
% Therefore
% \[
% \frac{\bigl|D_{\ell+1}\bigr|}{\bigl|D_{\ell}\bigr|}
% \;>\;
% 1
% \quad
% \text{for }\ell < 1.
% \]
% Consequently, higher-frequency modes (at $\ell+1$) do not decay faster than lower-frequency modes (at $\ell$), violating the standard Frequency Principle requirement that low frequencies be learned first.

% \noindent
% Thus, for $\ell=0$ or $\ell=1$, the update dynamics can \emph{break} the usual low-frequency dominance if the (non-vanishing) error terms push these early harmonics upward.  
% This completes the proof that the network can violate the frequency principle under the stated conditions.
\end{proof}

\end{document}
